
在長期互動的記憶系統中，同一事實會隨時間累積多個版本，辨識這些版本並採用其中的當前版本，此任務被稱為知識更新（knowledge update, KU），近期多個 benchmark 已將其視為記憶系統的核心能力之一。既有 KU 方法無論將判斷設計於寫入階段（write-time）或查詢階段（query-time），辨識哪些候選記憶對應到同一事實這一步皆由大型語言模型（LLM）的語意判斷負責，KU 的表現因而取決於 backbone 的判斷能力。然而實際部署常受隱私或成本的限制，可用的模型多為中小模型，此時 LLM 難以做出可靠的判斷；且當新事實與 LLM 既有知識衝突時，LLM 傾向以既有知識回答，而未依輸入的資訊做出更新。此外，將判斷設計於 write-time 的方法會在寫入當下就把判斷結果寫定於記憶，一旦誤判，正確版本在任何查詢之前便已從記憶中被移除。
 
基於上述條件，本文的目標是設計一個判斷不依賴 backbone 語意能力的 KU 判斷機制。本文提出將 KU 判斷延後至 query-time，並以事實的 triple 結構作為判斷依據，以其中的 $(\text{subject}, \text{predicate})$ 結構配對辨識同一事實，並以事實的寫入時序決定當前版本。此判斷為 deterministic，不涉及語意理解，因此不隨 backbone 能力變動，也不受 LLM 既有知識偏好的影響；LLM 僅在結構配對失效的少數候選上作為補救，不參與主要判斷，亦不修改記憶。
 
本文於兩個文本特性不同的 KU benchmark 上與既有方法比較。於 counterfactual 型 KU 情境下的 FC-SH，本方法四個 context length 的平均正確率為 92.5\%，高於最強 baseline 的 85.0\%，且於四個長度上維持在 91\% 至 94\%，不隨 context 變長而下滑；對 baseline 失敗題的逐題分析顯示，其錯誤幾乎都回退到與世界知識一致的舊值，此一致的失敗方向對應 LLM 對既有知識的偏好。於固定 context length 下改換不同判斷能力的 backbone，本方法相對最強 baseline 的差距隨 backbone 判斷能力減弱而擴大，隨著 backbone 能力足夠可靠而收斂。最後以 personal 型 KU 情境 LME-KU 作為對照，本方法相對 baseline 的差距明顯縮小，其中與本方法機制最接近的既有方法，由 counterfactual 情境下的領先 7.5 百分點反轉為落後 5.1 百分點（70.5\% 對 75.6\%）；此反轉與前述所預期的一致，因此佐證 counterfactual 情境上的差距確實來自 LLM 對既有知識的偏好。
 
整體而言，本文將主要的 KU 判斷從 LLM 移出，改以 deterministic 結構配對負責。於新進事實與 LLM 既有知識衝突且 backbone 判斷能力受限的情境下，讓 deterministic 機制負責主要判斷，LLM 僅作補救，是一種有效的記憶系統 KU 架構設計。

關鍵字：知識更新；記憶系統；知識衝突；確定性結構配對；大型語言模型